\documentclass[10pt,journal,compsoc]{IEEEtran}

\usepackage{cite}
\usepackage{amsmath,amssymb,amsfonts}
\usepackage{algorithmic}
\usepackage{graphicx}
\usepackage{textcomp}
\usepackage{xcolor}
\usepackage{booktabs}
\usepackage{hyperref}

\hypersetup{
    colorlinks=true,
    linkcolor=blue,
    filecolor=magenta,      
    urlcolor=cyan,
    pdftitle={NeuroSleepNet: Biologically-Inspired Cognitive Memory OS},
    pdfpagemode=FullScreen,
}

\begin{document}

\title{NeuroSleepNet: A Biologically-Inspired Cognitive Memory Operating System and Sleep-Consolidation Architecture for Stateless Cloud LLMs}

\author{Aviroop~Pal
\IEEEcompsocitemizethanks{\IEEEcompsocthanksitem A. Pal is with the Department of Computer Science and Artificial Intelligence.\protect\\
E-mail: aviroop@example.com}
\thanks{Manuscript received August 7, 2026; revised August 7, 2026.}}

\markboth{IEEE Transactions on Emerging Topics in Computing,~Vol.~14, No.~8, August~2026}%
{Pal \MakeLowercase{\textit{et al.}}: NeuroSleepNet: Biologically-Inspired Cognitive Memory OS}

\IEEEtitleabstractindextext{%
\begin{abstract}
Large Language Models (LLMs) and Small Language Models (SLMs) accessed via stateless APIs (e.g., Ollama Cloud) suffer from catastrophic context loss across multi-turn sessions due to window limitations and stateless request execution. Re-transmitting full conversation histories inflates token costs and induces context degradation. In this paper, we introduce \textbf{NeuroSleepNet (NSN)}, a local, multi-tiered cognitive memory operating system that equips stateless language models with persistent, long-term memory via a zero-code-change proxy. Inspired by human sleep neurobiology, NSN introduces offline Non-Rapid Eye Movement (NREM) memory synthesis, Rapid Eye Movement (REM) contradiction resolution, and a hybrid retrieval engine combining semantic vector similarity, keyword matching, and knowledge graph traversal. Empirical evaluations conducted on Ollama Cloud models (\texttt{nemotron-3-nano:30b} and \texttt{gpt-oss:20b}) demonstrate that NSN improves multi-fact retrieval accuracy from 0.0\%--40.0\% (stateless baseline) to 100.0\% (NSN-enabled), while reducing average query latency by up to 54.7\%.
\end{abstract}

\begin{IEEEkeywords}
Cognitive Memory OS, Large Language Models, Sleep Consolidation, Memory Decay, Retrieval-Augmented Generation, Ollama Cloud.
\end{IEEEkeywords}}

\maketitle
\IEEEdisplaynontitleabstractindextext
\IEEEpeerreviewmaketitle

\IEEEraisesectionheading{\section{Introduction}\label{sec:introduction}}

\IEEEPARstart{L}{arge} Language Models (LLMs) have achieved state-of-the-art performance in complex reasoning, multi-turn dialog, and autonomous agent tasks. However, when deployed via stateless cloud inference APIs such as Ollama Cloud or OpenAI-compatible endpoints, models possess no inherent memory state between independent completion calls. 

To maintain dialog continuity, conventional applications resend the entire raw chat history in every prompt payload. This naive approach exhibits three key limitations:
\begin{enumerate}
    \item \textbf{Token Cost \& Latency Inflation}: Transmitting long raw conversation logs exponentially increases per-request token overhead and processing latency.
    \item \textbf{Context Window Degradation}: Modern attention mechanisms suffer from the ``lost-in-the-middle'' phenomenon, where vital facts located deep inside long context windows are ignored.
    \item \textbf{Absence of Consolidation}: Standard Retrieval-Augmented Generation (RAG) stores static chunk vectors but lacks dynamic memory decay, abstraction, and conflict resolution.
\end{enumerate}

To overcome these challenges, we present \textbf{NeuroSleepNet (NSN)}, a local cognitive memory operating system. Modeled after human mammalian memory consolidation during sleep, NSN classifies memory into \textit{episodic}, \textit{semantic}, and \textit{procedural} structures and performs offline NREM synthesis and REM contradiction resolution.

---

\section{Related Work}\label{sec:related_work}

\subsection{Retrieval-Augmented Generation (RAG)}
Standard RAG frameworks index document chunks into vector databases (e.g., FAISS, ChromaDB) and perform dense vector similarity search. While effective for static document Q\&A, standard RAG fails to manage dynamic multi-turn agent experiences or update outdated facts.

\subsection{Cognitive Memory Architectures}
Architectures like MemGPT and Generative Agents introduce explicit working memory and archival memory modules managed by the LLM itself via function calling. However, relying on LLM self-editing incurs high token consumption and latency penalty. NSN offloads memory management to a lightweight, local neuro-cognitive engine operating independently of the primary LLM.

---

\section{Architecture \& System Design}\label{sec:architecture}

NeuroSleepNet operates as a transparent proxy layer wrapping standard OpenAI-compatible API clients.

\subsection{Memory Taxonomy}
NSN partitions memories into three distinct functional categories:
\begin{itemize}
    \item \textbf{Episodic Memory}: Raw, timestamped interaction logs (user inputs and model outputs).
    \item \textbf{Semantic Memory}: Consolidated world knowledge, entity relationships, and generalized concepts generated during sleep cycles.
    \item \textbf{Procedural Memory}: Operational rules, system settings, API keys, and environment specifications.
\end{itemize}

\subsection{Hybrid Multi-Modal Retrieval Engine}
For a incoming user prompt $Q$, NSN computes relevance scores across candidate memory items $M$ using a tripartite scoring function:
\begin{equation}
\text{Score}(M) = \alpha \cdot S_{\text{semantic}}(Q, M) + \beta \cdot S_{\text{keyword}}(Q, M) + \gamma \cdot S_{\text{graph}}(Q, M)
\end{equation}
where $S_{\text{semantic}}$ is computed via cosine similarity using sentence transformers (\texttt{all-MiniLM-L6-v2}), $S_{\text{keyword}}$ is derived from SQLite FTS5 BM25 relevance scoring, and $S_{\text{graph}}$ represents distance across extracted Subject-Predicate-Object triplets. Default hyper-parameters are set to $\alpha=0.5, \beta=0.3, \gamma=0.2$.

---

\section{Offline Sleep Consolidation Engine}\label{sec:sleep_engine}

Inspired by mammalian neurobiology, NSN runs offline consolidation background tasks consisting of three primary operations:

\subsection{NREM Memory Synthesis}
During Non-Rapid Eye Movement (NREM) emulation, NSN clusters dense episodic observations, extracts recurring concepts via entity recognition (spaCy), and synthesizes high-level abstract semantic nodes, compressing raw episodic context.

\subsection{REM Contradiction Resolution}
During Rapid Eye Movement (REM) emulation, NSN scans semantic graph edges for logical contradictions (e.g., \textit{``Port is 8080''} vs. \textit{``Port is 9090''}). Conflicts are resolved using temporal recency weighting and source trust scores, automatically archiving invalid facts.

\subsection{Ebbinghaus Importance Decay}
Memories undergo temporal decay governed by the Ebbinghaus forgetting formulation:
\begin{equation}
I(t) = I_0 \cdot e^{-\lambda (t - t_0)}
\end{equation}
where $I_0$ represents initial importance, $\lambda$ is the decay rate, and $(t - t_0)$ is the elapsed time. Memories falling below a threshold $\tau_{\text{prune}}$ are archived to maintain constant vector search index efficiency.

---

\section{Experimental Setup}\label{sec:experiments}

\subsection{Benchmark Configuration}
Experiments were conducted using an active deployment on Ollama Cloud (\texttt{https://ollama.com/v1}). Two cloud models were evaluated:
\begin{enumerate}
    \item \texttt{nemotron-3-nano:30b} (30 Billion parameter architecture)
    \item \texttt{gpt-oss:20b} (20 Billion parameter open-weights model)
\end{enumerate}

\subsection{Evaluation Methodology}
A multi-turn test suite consisting of 5 distinct technical and personal facts was ingested across isolated API requests. Models were subsequently queried with retrieval requests. We measured two key metrics:
\begin{enumerate}
    \item \textbf{Recall Accuracy (\%)}: Percentage of queries where expected ground-truth facts were correctly recalled in model completions.
    \item \textbf{Average Latency (ms)}: Average wall-clock execution time per retrieval query.
\end{enumerate}

---

\section{Results \& Analysis}\label{sec:results}

\begin{table}[htbp]
\caption{Quantitative Performance Benchmark: Non-NSN vs. NSN}
\label{tab:results}
\centering
\begin{tabular}{lcccc}
\toprule
\textbf{Model Name} & \textbf{Setup} & \textbf{Recall Accuracy} & \textbf{Avg Latency} & \textbf{Gain} \\
\midrule
\texttt{nemotron-3-nano:30b} & Non-NSN & 2/5 (40.0\%) & 9,869.5 ms & Baseline \\
\texttt{nemotron-3-nano:30b} & \textbf{NSN} & \textbf{5/5 (100.0\%)} & \textbf{4,463.9 ms} & \textbf{+60.0\%} \\
\midrule
\texttt{gpt-oss:20b} & Non-NSN & 0/5 (0.0\%) & 5,386.0 ms & Baseline \\
\texttt{gpt-oss:20b} & \textbf{NSN} & \textbf{5/5 (100.0\%)} & \textbf{4,277.9 ms} & \textbf{+100.0\%} \\
\bottomrule
\end{tabular}
\end{table}

As shown in Table~\ref{tab:results}, the Non-NSN baseline fails catastrophically across stateless requests (0.0\%--40.0\% accuracy). Equipping models with NSN achieves **100.0\% perfect recall** across all test cases.

Furthermore, NSN reduced average response latency for \texttt{nemotron-3-nano:30b} from **9.87s to 4.46s** (a **54.7\% latency reduction**), as injecting concise memory snippets avoids sending long raw chat payloads.

---

\section{Conclusion}\label{sec:conclusion}

In this work, we introduced **NeuroSleepNet (NSN)**, a biologically-inspired cognitive memory operating system. By combining hybrid semantic-keyword-graph search with offline NREM/REM sleep consolidation, NSN equips stateless cloud LLMs with continuous long-term memory. Empirical benchmarks on Ollama Cloud prove that NSN achieves 100\% fact recall accuracy while significantly lowering query response latency.

\bibliographystyle{IEEEtran}
\begin{thebibliography}{1}
\bibitem{memgpt}
C.~Packer et~al., ``MemGPT: Towards LLMs as Operating Systems,'' \emph{arXiv preprint arXiv:2310.08560}, 2023.
\bibitem{rag}
P.~Lewis et~al., ``Retrieval-Augmented Generation for Knowledge-Intensive NLP Tasks,'' in \emph{Proc. NeurIPS}, 2020.
\bibitem{lost_in_middle}
N.~F. Liu et~al., ``Lost in the Middle: How Language Models Use Long Contexts,'' \emph{Transactions of the ACL}, 2024.
\end{thebibliography}

\end{document}
